% Student assessment reconstructed from the complete matching solution source.
% Scoring ledger: Q1 C1=1.5+1.5, C2=1+1+1;
% Q2 same; Q3 C1=1.25+1.25, C2=0.75+1+0.75.
% Build twice with ../../../preamble/ in TEXINPUTS.
\makeatletter
\def\input@path{{../../../preamble/}}
\makeatother
\input{preamble_exams}
\graphicspath{{../../../preamble/}}
\usepackage{needspace}
\hypersetup{pdftitle={Grade 11 Term 2 Learning Evidence 1: C1 and C2},pdfauthor={Global Economics}}
\setYColumnWidth{0.36}
\linespread{1.04}
\raggedbottom
\newcolumntype{K}{>{\raggedright\arraybackslash}X}
\newcommand{\examtable}[2]{\par\noindent\begin{minipage}{\linewidth}\small
\renewcommand{\arraystretch}{1.08}\setlength{\tabcolsep}{3pt}
\begin{tabularx}{\linewidth}{#1}\toprule #2\bottomrule\end{tabularx}
\end{minipage}\par}
\let\ExamSubsectionBox\SubsectionBox
\RenewDocumentCommand{\SubsectionBox}{s m}{%
  \par\Needspace{4\baselineskip}% keep heading with opening prompt
  \begingroup
  \setlength{\linewidth}{\dimexpr\linewidth-2\fboxsep\relax}%
  \IfBooleanTF{#1}{\ExamSubsectionBox*{#2}}{\ExamSubsectionBox{#2}}%
  \endgroup\nobreak
}
% Shared headings use stepcounter, not refstepcounter: capture the automatic
% number explicitly for score-table references, without manual question IDs.
\makeatletter
\newcommand{\QuestionLabel}[1]{%
  \begingroup\edef\@currentlabel{\arabic{subsectionboxproblem}}%
  \label{#1}\endgroup
}
\makeatother
\newcommand{\cOneA}{Identify the population and both samples, with place, period, eligibility and sizes. Pair population parameters with sample statistics for the mean, standard deviation and proportion. Interpret every measure with units where appropriate; compute both sample proportions from their condition counts. Distinguish known quantities from those not established by the information.}
\newcommand{\cOneB}{Compare the second sample with the initial sample. Calculate changes in the mean, sample standard deviation and proportion, using percentage points for the last. Explain which population characteristics remain fixed and why sample differences do not establish changes in population parameters.}
\newcommand{\cTwoA}{Organize the initial study using the C2 reference categories: population parameters, sample statistics and procedure settings. Identify the initially unknown mean and its later established value. Calculate the standard error; identify confidence, significance, each tail area and critical value for every requested level.}
\newcommand{\cTwoB}{Using the initial sample, state the point estimate and calculate the margin of error, lower and upper endpoints, confidence interval and full width for each requested confidence level. Complete a reference-style calculation table; show formulas, substitutions, units and contextual interpretation. Use the known population standard deviation.}
\newcommand{\cTwoC}{Use only the subsequent studies to determine interval coverage at each level and its absolute difference from nominal confidence in percentage points. Interpret long-run coverage and finite-study variation. Distinguish interval coverage from sample means lying inside the initial interval, and avoid assigning probability to a fixed mean after an interval is calculated.}

\begin{document}
\noindent
\begin{tabularx}{\textwidth}{|C{2.8cm}|C{\dimexpr\textwidth-5.6cm-6\tabcolsep-4\arrayrulewidth\relax}|C{2.8cm}|}
\hline
\centering\vspace{3mm}\includegraphics[width=2.5cm]{logo.png}
& \CellCenter{\textbf{GLOBAL ECONOMICS}\par
\textbf{Grade 11 -- Learning Evidence 1}\par
\textbf{C1 and C2 Integrated Exam}\par
\textbf{Teacher: Nicolás López Cuéllar}}
& \CellCenter{\textbf{TERM 2}\par\textbf{2026--2027}}\\\hline
\end{tabularx}\par\smallskip
\LearningObjective{Distinguish fixed population parameters from variable sample statistics in economic and financial contexts; calculate known-population-standard-deviation confidence intervals and interpret their long-run coverage.}{%
\textbf{C1}: Distinguishes population parameters from sample statistics by interpreting and comparing population and sample means, standard deviations and proportions in economic and financial contexts.\par
\textbf{C2}: Compares point and interval estimation and interprets confidence through long-run coverage, distinguishing this from probability assigned to a particular calculated interval containing the parameter.}
\par\smallskip
\StudentInfo\par

\begin{multicols}{2}
\raggedcolumns
\setlength{\parskip}{3pt}
\SubsectionBox*{Criteria assessment}
Answer exactly three questions on \textbf{separate paper}; clearly identify each question and task. C1 has 8.5 available points and C2 has 8.5, for \textbf{17 points overall}.

\textbf{Each criterion requires at least 4 points from this exam to pass independently. Points cannot transfer between criteria.} An overall total alone does not establish a pass in both criteria.

Show reasoning, formulas, substitutions, units and contextual interpretations. Give exact probabilities where applicable. Use the supplied rounded critical values and retain all digits generated by them for endpoints and widths. Label any graphs you use; monetary axes must reflect actual monetary spacing. All scenario data are hypothetical.

\noindent\begin{minipage}{\linewidth}
\centering\small
\textbf{Available points}\par\smallskip
\begin{tabular}{|c|r|r|r|}
\hline\textbf{Question}&\textbf{C1}&\textbf{C2}&\textbf{Total}\\\hline
\ref*{q:1}&3&3&6\\\hline
\ref*{q:2}&3&3&6\\\hline
\ref*{q:3}&2.5&2.5&5\\\hline
\textbf{Total}&\textbf{8.5}&\textbf{8.5}&\textbf{17}\\\hline
\end{tabular}
\end{minipage}\par
\SubsectionBox*{Task reference material}
\begingroup\small
\textbf{C1 pairings.} Population $\mathcal P$, size $N$; sample $X$, size $n$; population/sample mean $(\mu,\bar x)$; standard deviation $(\sigma,s)$; proportion $(p,\hat p)$. A condition count is $k$, so $\hat p=k/n$. A sample-proportion change is $100(\hat p_B-\hat p_A)$ percentage points.

\textbf{C2 information table.} Use contiguous groups in this order: Population parameters ($\mu,\mu_{\mathrm{known}},\sigma$); Sample statistics ($\bar x,n,\mathrm{SE}$); Procedure settings ($C,\alpha,\alpha/2,z^*$). Columns are Symbol, Concept, Type and Meaning in context. As in the classroom reference, the derived SE is placed in the sample-statistics group; here it uses known $\sigma$ and fixed $n$, so it is not an estimate using $s$. The later-known notation refers to the same fixed mean, not a second parameter.

\textbf{Two-sided standard normal reference.} Use the following numerical critical values for all calculations. They are supplied rounded values; retain all digits generated by them when calculating endpoints and widths.
\examtable{KKK}{
\textbf{Confidence level}&\textbf{Area below positive critical value}&\textbf{Positive critical value}\\\midrule
90\%&0.950&1.645\\
95\%&0.975&1.960\\
99\%&0.995&2.576\\}
\[
\begin{aligned}
\alpha&=1-C, &\text{each tail}&=\alpha/2,\\
\mathrm{SE}&=\frac{\sigma}{\sqrt n}, &\mathrm{MoE}&=z^*\mathrm{SE},\\
L&=\bar x-\mathrm{MoE}, &U&=\bar x+\mathrm{MoE},\\
\mathrm{CI}&=[L,U], &W&=U-L=2\mathrm{MoE}.
\end{aligned}
\]
For each confidence level separately, a subsequent interval covers when
$L_i\leq\mu_{\mathrm{known}}\leq U_i$. Equality at an endpoint counts as coverage.
\[
\mathrm{Cov}(C)=\frac{N_{\mathrm{cover}}(C)}{N(C)}\,100\%.
\]
\[
E_{\mathrm{abs}}(C)=|\text{observed percentage}-100C|
\]
measured in percentage points.
The interval table uses Symbol, Concept, Calculation and Result and meaning in context, with separate endpoint rows. Proportions and confidence settings are unitless; a mean, standard deviation, standard error, margin, endpoint and width use the measurement's units.

\textbf{Common procedure assumptions.} The scenarios specify independent random studies, negligible within-study sampling fractions and sufficiently regular numerical populations for a normal approximation to the sampling distribution of the mean. The finite-population correction is negligible and omitted. Thus nominal normal coverage is an approximation, additionally subject to supplied critical-value rounding. Sample standard deviations describe samples; they do not replace known population standard deviations in this procedure.

\endgroup
\end{multicols}
\newpage
\begin{multicols}{2}
\raggedcolumns
\setlength{\parskip}{3pt}
\SubsectionBox{Retail electricity costs \score{C1-\sml{3}\quad C2-\sml{3}}}
\QuestionLabel{q:1}
\textbf{6 points overall.}

An energy-cost review concerns all 200,000 registered retail businesses operating in Medellin throughout October 2026. Each business contributes its total October electricity bill, measured in thousand COP. A separate yes-or-no condition records whether the business used electronic invoicing in October.

The initial independent random sample contains 100 businesses. Its mean bill is 240 thousand COP, its sample standard deviation is 58 thousand COP, and 40 businesses used electronic invoicing. A second independent random sample from exactly the same business population and October period also contains 100 businesses. Its mean bill is 252 thousand COP, its sample standard deviation is 66 thousand COP, and 55 businesses used electronic invoicing.

A certified spread-only population report makes the population standard deviation of October bills known as 60 thousand COP; it does not disclose the population mean. The two sample standard deviations remain separate descriptive quantities. The initial study estimates the initially unknown mean at 95\% confidence.

Assume the bill distribution is sufficiently regular, without dominant outliers, for a normal approximation to sample means of this size. Each study draws a simple random sample without replacement from the full fixed register. Samples are selected independently across studies; overlap between different samples is permitted. Each study selects only 0.05\% of the population, so dependence within a draw and the finite-population correction are negligible.

Ten subsequent independent studies each use 100 businesses, the same known population standard deviation and the same 95\% normal confidence procedure. These are additional studies, separate from both samples above. A later comprehensive investigation of the entire October population establishes its mean bill as 245 thousand COP. Repeated samples do not establish this value. The following subsequent intervals allow retrospective assessment; all endpoints are in thousand COP.

\examtable{KKK}{
\textbf{Subsequent study}&\textbf{Lower endpoint}&\textbf{Upper endpoint}\\\midrule
1&224.24&247.76\\
2&232.24&255.76\\
3&238.24&261.76\\
4&246.24&269.76\\
5&218.24&241.76\\
6&235.24&258.76\\
7&230.24&253.76\\
8&242.24&265.76\\
9&221.24&244.76\\
10&234.24&257.76\\
}

\columnbreak
\begin{enumerate}[label=(\alph*)]
\item \textbf{C1 Task 1} \sscore{C1-\sml{1.5}}\par
\cOneA
\item \textbf{C1 Task 2} \sscore{C1-\sml{1.5}}\par
\cOneB
\item \textbf{C2 Task 1} \sscore{C2-\sml{1}}\par
\cTwoA
\item \textbf{C2 Task 2} \sscore{C2-\sml{1}}\par
\cTwoB
\item \textbf{C2 Task 3} \sscore{C2-\sml{1}}\par
\cTwoC
\end{enumerate}

\end{multicols}
\newpage
\begin{multicols}{2}
\raggedcolumns
\setlength{\parskip}{3pt}
\SubsectionBox{Savings-account interest \score{C1-\sml{3}\quad C2-\sml{3}}}
\QuestionLabel{q:2}
\textbf{6 points overall.}

A bank reviews all 500,000 positive-balance personal savings accounts at its Cartagena branch network on October 31, 2026. The numerical measurement is the total interest credited to each eligible account during October 2026, in thousand COP. A separate yes-or-no condition records whether the account had mobile banking activated on October 31.

The initial independent random sample contains 144 eligible accounts, with mean October interest of 82 thousand COP, sample standard deviation of 33 thousand COP and 54 accounts with mobile banking activated. A second independent random sample of 144 accounts from exactly the same eligible population and period has mean interest of 79 thousand COP, sample standard deviation of 39 thousand COP and 72 accounts with mobile banking activated.

A certified spread-only report establishes the population standard deviation of October interest as 36 thousand COP while withholding the population mean. This known spread differs from both sample standard deviations. The initial study estimates the initially unknown population mean at 90\% and 99\% confidence.

Assume a sufficiently regular interest distribution without dominant outliers and a valid normal approximation for sample means of this size. Each study draws a simple random sample without replacement from the same full fixed register; draws across studies are independent and may overlap. Each sample contains only 0.0288\% of the population, so the finite-population correction is negligible.

Two hundred subsequent independent studies, separate from the two samples above, each select 144 accounts and use the same known population standard deviation and both confidence procedures. Each study builds both intervals around its own sample mean. The 99\% interval is wider and contains that study's 90\% interval. A later comprehensive investigation of all eligible accounts establishes mean October interest as 84 thousand COP; the repeated samples do not establish it. The table records mutually exclusive retrospective coverage categories for these same 200 studies.

\examtable{Kr}{
\textbf{Intervals covering the later established mean}&\textbf{Studies}\\\midrule
Both confidence levels&176\\
Only the 99\% level&20\\
Neither confidence level&4\\}

\columnbreak
\begin{enumerate}[label=(\alph*)]
\item \textbf{C1 Task 1} \sscore{C1-\sml{1.5}}\par
\cOneA
\item \textbf{C1 Task 2} \sscore{C1-\sml{1.5}}\par
\cOneB
\item \textbf{C2 Task 1} \sscore{C2-\sml{1}}\par
\cTwoA
\item \textbf{C2 Task 2} \sscore{C2-\sml{1}}\par
\cTwoB
\item \textbf{C2 Task 3} \sscore{C2-\sml{1}}\par
\cTwoC
\end{enumerate}

\end{multicols}
\newpage
\begin{multicols}{2}
\raggedcolumns
\setlength{\parskip}{3pt}
\SubsectionBox{Create an integrated scenario \score{C1-\sml{2.5}\quad C2-\sml{2.5}}}
\QuestionLabel{q:3}
\textbf{5 points overall.}
Invent a new economic or financial scenario integrating both criteria. Supply the shared contextual description and all data, then complete both C1 tasks and all three C2 tasks below for your own scenario. One confidence level is sufficient. Write the contextual description in ordinary language: numerals, percentages, written units and plain-language table headings are allowed, but reserve symbols, formulas, inequalities and interval notation for task reference material and solutions.

\textbf{Required construction checklist}
\begin{enumerate}
\item Define a population's place, fixed period, eligibility and size; define one numerical measurement with units and a separate binary condition.
\item Provide an initial independent random sample and a second independent random sample from exactly the same population and period. For each give size, mean, sample standard deviation and integer condition count between zero and that sample size. Ensure summaries are feasible for the numerical measurement.
\item State a known population standard deviation, distinguish it from both sample standard deviations, and explain that the population mean is initially unknown. Choose confidence and a supplied standard normal critical value.
\item Give suitable randomness, independence and normal-approximation assumptions. For finite sampling without replacement, keep the sampling fractions negligible. Subsequent studies must be independent draws from the same full fixed population and period.
\item Specify a number of subsequent studies, each with the initial sample size, known population standard deviation and confidence procedure. Supply all subsequent endpoints or complete coverage/miss counts. Endpoints must have the stated procedure's width; with multiple levels, evidence must respect nesting.
\item Include a later comprehensive population investigation establishing the mean of that same population for retrospective coverage. A clearly hypothetical investigation is acceptable; do not claim repeated samples establish the population mean.
\end{enumerate}

Scenario completeness and validity are assessed within the task marks below. Complete all five tasks for your own scenario.

\columnbreak
\begin{enumerate}[label=(\alph*)]
\item \textbf{C1 Task 1} \sscore{C1-\sml{1.25}}\par
\cOneA
\item \textbf{C1 Task 2} \sscore{C1-\sml{1.25}}\par
\cOneB
\item \textbf{C2 Task 1} \sscore{C2-\sml{0.75}}\par
\cTwoA
\item \textbf{C2 Task 2} \sscore{C2-\sml{1}}\par
\cTwoB
\item \textbf{C2 Task 3} \sscore{C2-\sml{0.75}}\par
\cTwoC
\end{enumerate}

\end{multicols}
\end{document}
